\chapter{The Roadmap: Twelve Phases}\label{ch:roadmap}

\begin{objectives}
\begin{itemize}
\item read the twelve-phase plan in \file{tasks/todo.md} as a dependency graph with a critical path;
\item say exactly what ``done'' meant for Phases 0--3 and what remains open inside them;
\item write and test the Phase 4 CMake change, the next step;
\item connect every later phase to the chapter that teaches its ideas.
\end{itemize}
{\small\color{muted}\textbf{Chapter references:} the repository files \file{tasks/todo.md} and \file{CHANGELOG.md}; Kerrisk~\cite{kerrisk2010tlpi} for the Linux-specific APIs.}
\end{objectives}

\section{The plan as a graph}
\begin{figure}[htbp]
\centering
\begin{tikzpicture}[
  ph/.style={draw,rounded corners=3pt,fill=codebg,text=fg,font=\footnotesize,align=center,minimum width=3.15cm,minimum height=1.25cm,inner sep=3pt},
  done/.style={ph,draw=grn,line width=1pt},
  next/.style={ph,draw=yel,line width=2pt,fill=yel!10!bg},
  todo/.style={ph,draw=acc},
  later/.style={ph,draw=pur,dashed},
  e/.style={-{Stealth[length=2mm]},draw=muted,thick}]
\node[done] (p0) at (0,0) {\textbf{P0} Risk\\historical VaR/CVaR};
\node[done] (p1) at (3.7,0) {\textbf{P1} Storage\\\texttt{ColumnStore} + V1 stub};
\node[done] (p2) at (7.4,0) {\textbf{P2} Ingestion\\\texttt{IngestionSource}};
\node[done] (p3) at (11.1,0) {\textbf{P3} Net\\interfaces + stubs};
\node[next] (p4) at (11.1,-2.0) {\textbf{P4} Net CMake\\Linux-only epoll guard};
\node[todo] (p5) at (7.4,-2.0) {\textbf{P5} v1 concurrency\\thread per connection};
\node[todo] (p6) at (3.7,-2.0) {\textbf{P6} v2 concurrency\\epoll (Linux)};
\node[todo] (p7) at (0,-2.0) {\textbf{P7} v3 (stretch)\\io\_uring};
\node[todo] (p8) at (7.4,-4.0) {\textbf{P8} Wire protocol\\binary codec};
\node[todo] (p9) at (3.7,-4.0) {\textbf{P9} Test coverage\\db + ingestion};
\node[todo] (p10) at (0,-4.0) {\textbf{P10} nanobind\\+ Python ingestion};
\node[later] (p11) at (0,-6.0) {\textbf{P11} gRPC facade\\(out of scope now)};
\draw[e] (p0)--(p1); \draw[e] (p1)--(p2); \draw[e] (p2)--(p3); \draw[e] (p3)--(p4);
\draw[e] (p4)--(p5); \draw[e] (p5)--(p6); \draw[e] (p6)--(p7);
\draw[e] (p5)--(p8); \draw[e] (p8)--(p9); \draw[e] (p9)--(p10); \draw[e] (p10)--(p11);
\draw[e] (p1) to[out=-60,in=60] (p9);
\node[font=\scriptsize,text=muted,anchor=west] at (5.0,-6.0) {\textcolor{grn}{\rule{8pt}{8pt}} done\quad \textcolor{yel}{\rule{8pt}{8pt}} next\quad \textcolor{acc}{\rule{8pt}{8pt}} planned\quad \textcolor{pur}{\rule{8pt}{8pt}} deferred};
\end{tikzpicture}
\caption{The twelve phases of \file{tasks/todo.md} as a dependency graph (TikZ). Edges are ``must precede''. Phase 7 is a
stretch goal off the main line; the main line to Phase 11 runs through the binary codec and the test pass.}
\label{fig:roadmap}
\end{figure}

\repomd[firstline=29,lastline=79]{tasks/todo.md}{The plan, verbatim (lines 29--79). Checked boxes are Phases 0--3. Some
items inside those phases remain open by design: the VaR test file, the on-disk layout, and the note about the Python
integration points.}{lst:todo}

\section{What ``done'' meant for Phases 0--3}
Decision D11 defined done narrowly. For Phase 0 it meant a real implementation (\cref{ch:risk}), spot-checked by hand,
with \file{tests/cpp/test_risk.cpp} still unchecked. For Phases 1--3 it meant headers with the final shape, \file{.cpp}
stubs that throw, and, for \file{db/}, inclusion in the build. According to the newest changelog entry, the existing suites
still passed afterwards: 16 of 16 C++ tests and 21 of 21 Python tests.

\begin{remarkbox}
The two kinds of ``done'' are visible in \cref{fig:arch}: \texttt{db/} has a solid blue border (built, stubbed),
\texttt{net/} a dashed yellow one (not built). Keeping \texttt{net/} out of the build was deliberate: adding it brings the
first platform-specific source file, and that deserves its own reviewed change.
\end{remarkbox}

\section{Phase 4: the next step}
Phase 4 adds the network sources to \cc{SCORE_CORE_CPP_SOURCES} and compiles the \cc{epoll} backend only on Linux.

\begin{code}[listing options={style=plaindark}]{src/CMakeLists.txt, Phase 4 sketch (not yet in the repository)}{lst:phase4}
list(APPEND SCORE_CORE_CPP_SOURCES
    net/json_codec.cpp
    net/connection.cpp
    net/server.cpp
)
if(CMAKE_SYSTEM_NAME STREQUAL "Linux")
    list(APPEND SCORE_CORE_CPP_SOURCES net/event_loop_epoll.cpp)
endif()
\end{code}

\begin{verified}
The snippet was configured in a scratch CMake project (CMake from Homebrew): on macOS it printed
\texttt{system=Darwin} and the three network sources; configured with \texttt{-DCMAKE\_SYSTEM\_NAME=Linux} it printed
\texttt{system=Linux} and added \texttt{net/event\_loop\_epoll.cpp}.
\end{verified}

Two engineering points matter more than the syntax.
\begin{enumerate}[label=\textbf{\color{yel}\arabic*.}]
\item \textbf{Guard the build and the source.} The CMake guard keeps the file out of non-Linux builds; an
{\ttfamily\color{org}\#if defined(\_\_linux\_\_)} inside the file (as the plan for Phase 6 says) keeps it harmless if someone adds it by hand.
Platform tests by feature (``does \file{sys/epoll.h} exist?'') are more robust than by name; \cref{ex:feature} explores this.
\item \textbf{Prove both platforms build.} The plan's second Phase 4 item is to confirm \cc{score_core} builds on
macOS and Linux. Until a CI job runs on both, D4's promise is untested.
\end{enumerate}

\section{The phases and the chapters that teach them}
\begin{center}
\footnotesize
\begin{tabularx}{\linewidth}{@{}>{\leavevmode\bfseries\color{yel}}p{0.8cm} >{\raggedright\arraybackslash}X >{\raggedright\arraybackslash}p{3.0cm} >{\raggedright\arraybackslash}p{4.4cm}@{}}
\toprule
\textcolor{muted}{Phase} & \textcolor{muted}{Deliverable} & \textcolor{muted}{Status} & \textcolor{muted}{Where taught}\\
\midrule
0 & historical VaR/CVaR & implemented; test file open & \cref{ch:risk,ch:float}\\
1 & \cc{ColumnStore}, \cc{ColumnStoreV1} stub & scaffolded, built & \cref{ch:columnar}\\
2 & \cc{IngestionSource} & scaffolded & \cref{sec:gbm}\\
3 & \cc{EventLoop}, \cc{FrameCodec}, \cc{JsonFrameCodec}, \cc{Connection}, \cc{Server} & scaffolded, not built & \cref{ch:tcp,ch:conc,ch:framing}\\
4 & net CMake wiring, Linux-only epoll guard & \textcolor{yel}{next} & this chapter\\
5 & \cc{BlockingThreadEventLoop} + loopback echo test & planned & \cref{sec:tpc}, \cref{ex:blocking}\\
6 & \cc{EpollEventLoop} (Linux) & planned & \cref{thm:epoll}, \cref{ex:et}\\
7 & \cc{IoUringEventLoop} (stretch) & planned & \cref{fig:uring}, \cref{ex:uringadapt}\\
8 & \cc{BinaryFrameCodec} & planned & \cref{def:varint}, \cref{ex:binary}\\
9 & tests for store and ingestion & planned & \cref{thm:scan}, \cref{ex:scanedge}\\
10 & nanobind exposure; Python ingestion wiring & planned & \cref{ex:yf}, \cref{ex:gil}\\
11 & gRPC facade & deferred & \cref{ex:grpc}\\
\bottomrule
\end{tabularx}
\end{center}

\begin{exercises}
\ex{ex:phase4}{\lvA} The three \file{src/net/*.cpp} files already pass \cc{clang++ -std=c++23 -fsyntax-only} on macOS
(checked for this tutorial). Apply \cref{lst:phase4} as written and configure on Linux: what does CMake report, given that
\file{net/event_loop_epoll.cpp} does not exist yet? Choose between adding a stub file in Phase 4 and deferring the
\cc{if} block to Phase 6, and justify the choice against the plan's ``reviewed, isolated CMake change''.

\ex{ex:feature}{\lvB} Replace the name test \cc{CMAKE_SYSTEM_NAME STREQUAL "Linux"} with a feature test using
\cc{check_include_file(sys/epoll.h ...)} or \cc{check_symbol_exists(epoll_create1 ...)}. Name one platform where the two
tests disagree, and say which behaviour you want there.

\ex{ex:ci}{\lvB} Design a CI matrix that keeps D4 honest: which jobs, on which runners, building which backends, running
which tests. What is the cheapest matrix that would have caught a {\ttfamily\color{org}\#include <sys/epoll.h>} leaking into \cc{server.cpp}?

\ex{ex:order8}{\lvB} The graph puts the binary codec (P8) after v1 concurrency (P5) and in parallel with epoll (P6).
Argue for or against doing P8 before P5. Which tests written for P5 would then also test P8?

\ex{ex:gil}{\lvC} Phase 10 asks for ``the Python surface for \cc{net::Server}''. A reactor thread that calls back into
Python must hold the GIL; one that never does can release it. Design the binding so that \cc{Server::run} releases the
GIL, handlers run in C++, and Python can still stop the server. What happens to a Python exception raised in a callback?

\ex{ex:grpc}{\lvC} Phase 11 evaluates gRPC ``as a facade over the \cc{FrameCodec}/\cc{Server} abstraction''. gRPC runs
over HTTP/2 with its own framing. Explain why gRPC cannot simply be another \cc{FrameCodec}, and sketch what ``facade''
should mean instead.
\end{exercises}
